\documentclass[11pt]{article}
\usepackage[margin=24mm]{geometry}
\usepackage[T1]{fontenc}
\usepackage{amsmath,amssymb,amsthm,mathtools,mathrsfs,longtable,booktabs,array,graphicx}
\usepackage[hidelinks]{hyperref}
\usepackage{xurl}
\hypersetup{hypertexnames=false}
\allowdisplaybreaks
\setlength{\emergencystretch}{3em}
\title{Finite native moments and the original arithmetic observation}
\author{}\date{20 September 2026}
\begin{document}\maketitle
Both complete proofs below retain the actual arithmetic measure, its mass,
all complex phases, the original quotient fibres and the full conductor.
Their human sources are cited at the point of use. These formulas provide
finite moment certificates, not assigned numerical values of those moments
or an evaluated growing-degree projected current.

% BEGIN COMPLETE NATIVE FINITE-MOMENT PROOFS026

\clearpage\begingroup
\part*{One scalar resolvent uncertainty in the original observation}


The matrix Gauss/Gauss--Radau method and its monotone two-sided bounds
are established human mathematics. Zimmerling, Druskin and Simoncini
\cite{rrnative:ZDS}, Sections 3--5, give their block continued-fraction formulation,
with explicit attribution to the earlier scalar quadrature literature.
Their original author LaTeX, both appendices and bibliography were read
in full for this calculation. Here the original arithmetic polynomial
block has exact linear dependences. We prove its complete finite receiver
directly, including the mass and complex coordinate factors. We do not
import a nondeflation assumption, the paper's scalar potential-theory
asymptotics, or its numerical observations at imaginary shifts.

\section{The unchanged arithmetic measures}
Keep the complete quartet divisor $h$, $g=2\xi$, and
\[
 w_h(y)=\frac{|(g/h)(1/2+iy)|^2}{2\pi},\qquad
 m_k=w_h^{*k},\qquad c=k/2,
\]
with the original even convolution measure and mass $\mu_h^k$.
The parity receiver \cite{rrnative:PP} proves the exact measures
\[
 d\nu_0(x)=m_k(\sqrt x)x^{-1/2}dx,\quad d\nu_1=x\,d\nu_0,
 \quad s(x)=\frac{\tanh(\pi\sqrt x/2)}{\sqrt x},\quad
 d\rho_b=s\,d\nu_b\quad(b=0,1).
\]
The value $s(0)$ is $\pi/2$. The density is positive almost everywhere
on $(0,\infty)$: the entire nonzero $g/h$ has at most isolated zeros
on the vertical line, and convolution of its positive-almost-everywhere
density preserves that property. The previously proved exponential
bound in $|y|$ implies that every polynomial moment below is finite.
For each of the two measures separately set
\[
 \mu_j^{(b)}=\int_0^\infty x^j\,d\rho_b(x).
\]
These are the actual tilted native moments. They are not the moments
of the raw Pollaczek weight, and $\mu_0^{(b)}$ is not replaced by one.
In the next sections write $\rho$ for either one, retaining that choice.

For $v_n(x)=(1,x,\ldots,x^n)^T$, the finite resolvent observation is
\begin{equation}\tag{RR1}\label{rrnative:eq:F}
 F_n(a)=\int\frac{v_n(x)v_n(x)^*}{x+a}\,d\rho(x),\qquad a>0.
\end{equation}
Multiplication $Xf=xf$ has its original domain in $L^2(\rho)$.
The map $B_nz=v_n^Tz$ satisfies
$B_n^*B_n=[\mu_{i+j}]_{i,j=0}^n$, not $I$.
The exact block Krylov space is
\begin{equation}\tag{RR2}\label{rrnative:eq:deflation}
 \operatorname{span}\{X^jB_nz:0\le j<m\}
       =\mathcal P_{n+m-1},\qquad\dim=n+m.
\end{equation}
Indeed its generators are exactly the powers of degrees from zero
to $n+m-1$; positivity of the density makes those powers independent
in $L^2(\rho)$. For $n>0$ and $m>1$ this is smaller than $m(n+1)$.
Thus the following calculation uses the actual polynomial space,
without pretending that this block is nondeflating.

\section{Two original-moment matrices and their exact errors}
Choose any integer $r\ge n+1$. All matrices in this section are in
the original monomial coordinates. Put
\[
 H_r=[\mu_{i+j}]_{i,j=0}^{r-1},\quad
 D_r=[\mu_{i+j+1}]_{i,j=0}^{r-1},\quad
 C_{r,n}=[\mu_{i+j}]_{\substack{0\le i<r\\0\le j\le n}},
 \quad e=(0,\ldots,0,1)^T.
\]
Both $H_r,D_r$ are positive definite, since their forms are integrals
of squared nonzero polynomials against $\rho$ and $x\rho$.
Define
\begin{equation}\tag{RR3}\label{rrnative:eq:forms}
 \kappa_r=(e^*D_r^{-1}e)^{-1},\quad
 \widetilde D_r=D_r-\kappa_ree^*,\quad R_r(a)=(D_r+aH_r)^{-1},
\end{equation}
and
\begin{equation}\tag{RR4}\label{rrnative:eq:bounds}
 F^-_{n,r}=C_{r,n}^*R_r(a)C_{r,n},\qquad
 F^+_{n,r}=C_{r,n}^*(\widetilde D_r+aH_r)^{-1}C_{r,n}.
\end{equation}
The matrix $\widetilde D_r$ is positive semidefinite with a
one-dimensional nullspace. To prove this, minimize $q^*D_rq$ under
$e^*q=t$. The attained minimum is $\kappa_r|t|^2$, at
$q=t\kappa_rD_r^{-1}e$. Subtracting that minimum proves the
assertion, including its kernel. Consequently both inverses in RR4
exist for every actual $a>0$.

Let $p_r$ be the monic degree-$r$ polynomial orthogonal for $\rho$,
and let $t_{r-1}$ be the monic degree-$(r-1)$ polynomial orthogonal
for $x\rho$. Their existence and uniqueness follow by solving the
positive Gram systems. The preceding minimum also proves
\begin{equation}\tag{RR5}\label{rrnative:eq:kappa}
 \int x|t_{r-1}(x)|^2d\rho(x)=\kappa_r.
\end{equation}
Neither polynomial vanishes at a negative real point. For example,
if $p_r(-a)=0$, then $q=p_r/(x+a)$ is a polynomial of degree $r-1$
and orthogonality would give
$0=\int p_rq\,d\rho=\int(x+a)q^2d\rho>0$.
The identical argument against $x\rho$ proves $t_{r-1}(-a)\ne0$;
for $r=1$, $t_0=1$.

For any complex $f\in\mathcal P_n$, let $q$ be the polynomial with
coefficient vector $R_rC_{r,n}[f]$. Its defining equations give
\[
 \int\overline u\{f-(x+a)q\}\,d\rho=0\qquad(u\in\mathcal P_{r-1}).
\]
The residual has degree at most $r$ and leading coefficient $-q_{r-1}$.
Therefore
\[
 f-(x+a)q=-q_{r-1}p_r,
 \qquad q_{r-1}=-\frac{f(-a)}{p_r(-a)}.
\]
Expanding the complete square, with these orthogonality equations,
now proves the exact Gauss error:
\begin{equation}\tag{RR6}\label{rrnative:eq:GaussError}
 [f]^*(F_n-F^-_{n,r})[f]
 =\int\frac{|f-(x+a)q|^2}{x+a}\,d\rho
 =\frac{|f(-a)|^2}{p_r(-a)^2}
              \int\frac{p_r(x)^2}{x+a}\,d\rho(x).
\end{equation}

For the other inverse in RR4, write its polynomial as $\widetilde q$.
Its equations read
\[
 \int\overline u\{f-(x+a)\widetilde q\}\,d\rho
       =-\kappa_r\overline{u_{r-1}}\widetilde q_{r-1}.
\]
Orthogonality for $x\rho$ and RR5 show that the polynomial
$f-(x+a)\widetilde q+\widetilde q_{r-1}xt_{r-1}$ is orthogonal
to every degree-$r-1$ polynomial and has that degree or less.
It is zero. Evaluation at $-a$ therefore gives
\[
 f-(x+a)\widetilde q=-\widetilde q_{r-1}xt_{r-1},
 \qquad\widetilde q_{r-1}=\frac{f(-a)}{a\,t_{r-1}(-a)}.
\]
In $\int \overline f\,xt_{r-1}/(x+a)d\rho$, divide
$\overline f(x)-\overline f(-a)$ by $x+a$. Its degree is at
most $r-2$, so that polynomial contribution vanishes by
orthogonality for $x\rho$. Applying the same division to
$t_{r-1}(x)-t_{r-1}(-a)$ gives the exact Radau error
\begin{equation}\tag{RR7}\label{rrnative:eq:RadauError}
 [f]^*(F^+_{n,r}-F_n)[f]
 =\frac{|f(-a)|^2}{a\,t_{r-1}(-a)^2}
               \int\frac{x\,t_{r-1}(x)^2}{x+a}\,d\rho(x).
\end{equation}
For constant polynomials the divided difference is zero; the proof
therefore includes $r=1$. RR6--7 prove the two-sided matrix bounds
without a finite-dimensional ambient operator or a spectral gap at zero.

\section{The entire uncertainty is one scalar, with its phase retained}
Put $z_n(a)=v_n(-a)$. Polarization of RR6--7 proves
\begin{equation}\tag{RR8}\label{rrnative:eq:rankone}
 F_n(a)=F^-_{n,r}(a)+\lambda_r(a)z_n(a)z_n(a)^*,
 \qquad 0\le\lambda_r(a)\le\gamma_r(a),
\end{equation}
where the exact, still native, scalar is
\[
 \lambda_r(a)=\frac1{p_r(-a)^2}
                 \int\frac{p_r(x)^2}{x+a}\,d\rho(x).
\]
The upper scalar needs only the displayed finite moments:
\begin{equation}\tag{RR9}\label{rrnative:eq:gamma}
 \gamma_r(a)=\frac{\kappa_r}
 {p_r(-a)^2\{1-\kappa_r e^*R_r(a)e\}}.
\end{equation}
To prove RR9, the rank-one inverse identity gives
\[
 (\widetilde D_r+aH_r)^{-1}-R_r
 =\frac{\kappa_r}{1-\kappa_r e^*R_re}R_ree^*R_r.
\]
Its denominator is positive: conjugate the positive definite
$\widetilde D_r+aH_r$ by the positive square root of $R_r$.
The already proved highest-coefficient identity gives
$e^*R_rC_{r,n}=-z_n(a)^*/p_r(-a)$.
Substitution proves $F^+-F^-=\gamma_rz_nz_n^*$, including the
complete mass. Indeed multiplying $\rho$ by any positive $M$
multiplies $H,D,C,\kappa,F^\pm,\lambda,\gamma$ by $M$,
while leaving both monic polynomials unchanged.

There is also a coordinate-level proof that only one scalar can occur.
For $i+j=d\ge1$, polynomial division proves
\begin{equation}\tag{RR10}\label{rrnative:eq:division}
 (F_n(a))_{ij}
 =\sum_{l=0}^{d-1}(-a)^l\mu_{d-1-l}+(-a)^d(F_n(a))_{00}.
\end{equation}
The case $d=0$ consists of the scalar itself. RR8 supplies certified
bounds for that exact scalar, rather than separate unrelated entry bounds.
In particular, for any specified complex coefficient vectors $u,v$,
\begin{equation}\tag{RR11}\label{rrnative:eq:phase}
 u^*(F_n-F^-)v=\lambda_r(a)(u^*z_n(a))(z_n(a)^*v).
\end{equation}
Thus it lies on one specified complex
line segment, not in an independently chosen error disc. Equivalently,
if $f_u=v_n^Tu$ and $f_v=v_n^Tv$, the factor is
$\lambda_r(a)\overline{f_u(-a)}f_v(-a)$.

\section{Positive resolvent sums and the original receiving maps}
The parity construction uses
$B_n(a)=\int vv^*x/(x+a)d\rho=H_{n+1}-aF_n(a)$.
Consequently define its moment-computable bounds
\begin{equation}\tag{RR12}\label{rrnative:eq:B}
 B^-_{n,r}=H_{n+1}-aF^+_{n,r},\quad
 B^+_{n,r}=H_{n+1}-aF^-_{n,r},\quad
 B^-_{n,r}\preceq B_n\preceq B^+_{n,r}.
\end{equation}
The lower matrix is positive semidefinite. One direct proof is to
use the previous Radau construction against the polynomials of
degree at most $r-2$ for the measure $x\rho$: its complementary
minimum is positive. Alternatively, since
$\widetilde D_r\succeq0$, the inverse inequality
$(\widetilde D_r+aH_r)^{-1}\preceq a^{-1}H_r^{-1}$ gives
$aF^+_{n,r}\preceq C_{r,n}^*H_r^{-1}C_{r,n}=H_{n+1}$.
This second proof also covers $r=1$.

At each original positive shift $a_j=4j^2$, independently choose
an actual degree $r_j\ge n+1$. The finite source observation
\[
 H^{[J]}=\frac2\pi H_{n+1}
          +\frac4\pi\sum_{j=1}^JB_n(4j^2)
\]
has lower and upper matrices $H^{[J],-},H^{[J],+}$ obtained by
using $B^-_{n,r_j},B^+_{n,r_j}$. Its exact uncertainty is
\begin{equation}\tag{RR13}\label{rrnative:eq:zonotope}
 H^{[J]}=H^{[J],+}
    -\frac4\pi\sum_{j=1}^J a_j\lambda_{r_j}(a_j)z_n(a_j)z_n(a_j)^*.
\end{equation}
Every coefficient retains its original interval $[0,\gamma_{r_j}(a_j)]$.
This is an enclosing real-parameter matrix family; independence of
those coefficients is not asserted. Each is generated by the same
unchanged native measure. There is no presumed sign of a complex
cross product.

Form the direct sum for the actual even and odd degrees, retaining
their separate measures, masses and shifts. In the original complex
$S$ coordinates the source Gram is $L_N^*HL_N$, where
\[
 (T_N)_{lj}=\binom jl c^{j-l}i^l\ (l\le j),\qquad
 L_N=\Pi_NT_N.
\]
The permutation $\Pi_N$ orders even powers first. RR13 is transported
by this literal congruence; each column $z$ becomes $L_N^*z$.
No unit modulus factors or translation terms are discarded.

The earlier exact native tail \cite{rrnative:TR} gives
$0\preceq H-H^{[J]}\preceq(2J)^{-1}X_H$.
Its original-moment alternative is
$X_H=[\int x^{i+j+1}d\nu_b]$ in each parity block. Thus
\begin{equation}\tag{RR14}\label{rrnative:eq:source}
 H^{[J],-}\preceq H\preceq H^{[J],+}+(2J)^{-1}X_H.
\end{equation}
The lower matrix is positive definite because its first summand
is $(2/\pi)H_{n+1}$ for each actual tilted measure.

For the unchanged onto observation $A$ in these parity coordinates,
the attained metric is $Q_A(G)=(AG^{-1}A^*)^{-1}$: its minimizing
representative is $G^{-1}A^*Q_A(G)u$. Subtract it from any other
representative in the same fibre and expand the square to prove
orthogonality to $\ker A$, hence the formula. Minimization in the
same fibres proves from RR14
\begin{equation}\tag{RR15}\label{rrnative:eq:quotient}
 Q_A(H^{[J],-})\preceq Q_A(H)
       \preceq Q_A(H^{[J],+}+(2J)^{-1}X_H).
\end{equation}
The full conductor has observation $\bar CA$, so the identical
formula with that actual onto map proves its two-sided metric
bound. This includes its entire kernel, not only a four-dimensional
chosen restriction. Restriction afterwards to the original signed
state map is an ordinary congruence and preserves the bound.

Write these endpoint bounds as $Q_i^-\preceq Q_i\preceq Q_i^+$.
Their positive determinants give, for the original signed return,
\begin{equation}\tag{RR16}\label{rrnative:eq:endpoints}
 \log\frac{\det Q_1^-\det Q_2^-}{\det Q_3^+\det Q_4^+}
 \le\log\frac{\det Q_1\det Q_2}{\det Q_3\det Q_4}
 \le\log\frac{\det Q_1^+\det Q_2^+}{\det Q_3^-\det Q_4^-}.
\end{equation}
For example determinant monotonicity follows by conjugating the
two positive matrices and multiplying eigenvalues at least one;
the original frames and their determinants are not replaced.
All quantities in these bounds use finite native moments, apart
from the explicitly displayed native tail if one elects not to
evaluate its next moments. The equations do not supply numerical
values of those moments, uniform convergence at growing packet
degree, a sign for the projected-current asymptotic, or an RH conclusion.

\begin{thebibliography}{9}
\bibitem{rrnative:ZDS} J. Zimmerling, V. Druskin and V. Simoncini,
\emph{Monotonicity, bounds and acceleration of Block Gauss and
Gauss--Radau quadrature for computing $B^T\phi(A)B$}, original
author source of arXiv:2407.21505v3 (9 January 2025),
Sections 3--5, especially source labels
\texttt{lemma:T\_LDLt}, \texttt{prop:Pencil2Stieltjes},
\texttt{lem01} and \texttt{prop:03}.
\url{https://arxiv.org/abs/2407.21505v3}.
The arXiv abstract page retains ``extrapolation'' in the title;
the title above transcribes the actual v3 author TeX.
The earlier scalar quadrature works cited there are human
provenance leads, not claimed here as independently read sources.
\bibitem{rrnative:PP} Split-Zero programme, \emph{Exact native parity transport
into the Pollaczek pair and certified finite source--observation bounds},
complete proof \texttt{PARITY\_POLLACZEK\_NATIVE\_RECEIVER.tex},
SZ-20260920-019, with the original Aptekarev--L\'opez
Lagomasino--Mart\'{\i}nez--Finkelshtein weight-ratio citation.
\url{https://github.com/KokunoYumeto/zeta-function-research-reader/tree/ab45e221579a0d48ce2885b4ecdf1a6aefc24a20/workbenches/splitzero-tandem/continuations/20260920-native-parity-conductor-return}.
\bibitem{rrnative:TR} Split-Zero programme, \emph{Native resolvent truncation at
growing degree and the original conductor minimum}, same complete
source edition, \texttt{NATIVE\_TRUNCATION\_CONDUCTOR\_RETURN.tex}, TR1--10.
\end{thebibliography}

\endgroup

\clearpage\begingroup
\part*{Complete finite native moment enclosures and original endpoint transport}
\newtheorem{mmnativetheorem}{Theorem}
\newtheorem{mmnativelemma}[mmnativetheorem]{Lemma}
\newtheorem{mmnativeproposition}[mmnativetheorem]{Proposition}
\newcommand{\mmnativeC}{\mathbb C}
\newcommand{\mmnativeR}{\mathbb R}
\newcommand{\mmnativediag}{\operatorname{diag}}
\newcommand{\mmnativeim}{\operatorname{im}}
\newcommand{\mmnativetr}{\operatorname{tr}}
\newcommand{\mmnativenorm}[1]{\left\lVert#1\right\rVert}


This note closes the finite-resolvent step left in the complete native
parity/Pollaczek calculation \cite{mmnative:PP} and its conductor continuation
\cite{mmnative:TR}.  Every resolvent below is taken against an actual native
measure.  Finite moment matrices provide two-sided bounds, exact
variational error identities, and computable positive-semidefinite gap
certificates.  The unknown native moments remain explicit inputs; no
numerical value, phase estimate, or substitute measure is assigned to
them.  The original observation and conductor kernels and the signs of
the four endpoint errors are retained.

\section{Unchanged objects, parity measures, and the unresolved integrals}

Keep the full quartet divisor and original arithmetic source
\begin{equation}\label{mmnative:eq:native}
 \begin{gathered}
 \rho_*=\tfrac12+\delta+i\gamma,\quad0<\delta<\tfrac12,\quad\gamma>2,
 \quad m\ge1,\quad g=2\xi,\\
 h(s)=\prod_{z\in\{\rho_*,\bar\rho_*,1-\rho_*,1-\bar\rho_*\}}(s-z)^m,
 \quad w_h(y)=\frac{|(g/h)(1/2+iy)|^2}{2\pi},\\
 m_k=w_h^{*k},\quad k=4\ell+1\ge9,\quad
 \mu_h=\int_\mmnativeR w_h(y)\,dy,\quad S=c+iy,\quad c=k/2,\\
 e_k=1+k(m-1),\quad q=e_k(k+1)^2,\quad
 \chi(S)=\prod_{a,b=0}^k
 [S-c-(2a-k)\delta-i(2b-k)\gamma]^{e_k}.
 \end{gathered}
\end{equation}
The full multiplicity $m$ and the physical Taylor unit are unchanged.
This is the stipulated quartet domain, not an assertion of existence of
an off-line zero.  The complete native providers \cite{mmnative:CAI,mmnative:HG} prove
evenness, positivity, and the actual envelope
\begin{equation}\label{mmnative:eq:envelope}
 0<a_k(1+y^2)^{-M}\rho_\sigma(y)\le m_k(y)\le A_k\rho_\sigma(y),
 \quad M=21+4m,\quad
 \rho_\sigma(y)=\frac{|\Gamma(1/4+iy/2)|^2}{2\pi},
\end{equation}
where the constants retain their original definitions
\begin{equation}\label{mmnative:eq:constants}
 \begin{gathered}
 \alpha=\pi/2,\quad p=8m-9/2,\quad B_h=42+8m,\\
 c_\Gamma=\sqrt2e^{-7/6},\quad
 C_\Gamma=\sqrt{2\sqrt5}e^{2/3},\quad
 \rho_\sigma(y)\le C_\Gamma e^{-\alpha|y|}(1+|y|)^{-1/2},\\
 \vartheta_h=\int_{-1}^1w_h(y)\,dy,\quad
 M_\alpha=\int_\mmnativeR e^{\alpha y}w_h(y)\,dy,\\
 a_k=\frac{c_h\vartheta_h^{k-3}e^{-\alpha(k-3)}(k-2)^{-B_h}}
 {C_\Gamma2^{B_h/2}},\qquad
 A_k=\frac{C_h^{\rm tilt}}{c_\Gamma}k^{p+1}M_\alpha^{k-1}.
 \end{gathered}
\end{equation}
Here $c_h$ and $C_h^{\rm tilt}$ are precisely the complete CAI
provider's TW7 and BT12 constants, not newly postulated numbers.

Put $\zeta_r=\int y^rm_k(y)\,dy$ and retain both parity measures
\begin{equation}\label{mmnative:eq:parities}
 d\nu_0(x)=\frac{m_k(\sqrt x)}{\sqrt x}\,dx,\qquad
 d\nu_1(x)=x\,d\nu_0(x),\qquad x>0.
\end{equation}
For the original $f(y)=E(y^2)+yO(y^2)$, evenness and the two
half-axis substitutions give
\begin{equation}\label{mmnative:eq:paritynorm}
 \int_\mmnativeR|f(y)|^2m_k(y)dy
 =\int|E|^2d\nu_0+\int|O|^2d\nu_1.
\end{equation}
The masses are exactly
\[
 \nu_0((0,\infty))=\mu_h^k,\qquad
 \nu_1((0,\infty))=\zeta_2=k\eta_2\mu_h^{k-1},
 \quad\eta_2=\int y^2w_h(y)dy.
\]
The second identity follows by the convolution multinomial formula;
the mixed first moments vanish by evenness.  No mass is replaced by one.

The unscaled Pollaczek ratio supplied by Aptekarev, L\'opez Lagomasino
and Mart\'{\i}nez--Finkelshtein \cite{mmnative:ALM}, and the reciprocal identity
proved by the companion's Green-kernel calculation \cite{mmnative:PP}, are
\begin{equation}\label{mmnative:eq:s}
 \begin{gathered}
 s(x)=\frac{\tanh(\alpha\sqrt x)}{\sqrt x}
 =\frac4\pi\sum_{r=0}^\infty\frac1{x+(2r+1)^2},
 \quad0<s(x)\le\alpha,\\
 \frac1{s(x)}=\frac2\pi+\frac4\pi
       \sum_{j=1}^\infty\frac{x}{x+4j^2}.
 \end{gathered}
\end{equation}
These identities are used with the actual tilted measures
\begin{equation}\label{mmnative:eq:tilts}
 d\widetilde\nu_\epsilon=s(x)d\nu_\epsilon,
 \qquad\epsilon\in\{0,1\},
 \qquad\widetilde\nu_1=x\widetilde\nu_0.
\end{equation}
Their moments and masses are not the raw Pollaczek moments.  For a
parity degree $n$ the unresolved matrix in \cite{mmnative:TR} is exactly
\begin{equation}\label{mmnative:eq:target}
 B_{\epsilon,n}(a)=\int_0^\infty
 \frac{x}{x+a}v_n(x)v_n(x)^*\,d\widetilde\nu_\epsilon(x),
 \quad v_n=(1,x,\ldots,x^n)^T,\quad a=4j^2.
\end{equation}
The following construction evaluates certified intervals for this
matrix using a finite sequence of its actual measure's moments.

\section{Two complementary finite moment principles}

In this section $d\lambda$ is one of the actual measures
$d\nu_\epsilon$ or $d\widetilde\nu_\epsilon$.  Write
$\mu_r=\int x^r\,d\lambda$.  For $\ell\ge0$ set
\[
 H_t^{(\ell)}=[\mu_{i+j+\ell}]_{i,j=0}^t,
 \qquad H_t=H_t^{(0)}.
\]
For a target degree $n$ and an auxiliary degree $t\ge n$, put
\begin{equation}\label{mmnative:eq:rectangular}
 C_t=[\mu_{i+j}]_{\substack{0\le i\le t\\0\le j\le n}},
 \quad D_t=[\mu_{i+j+1}]_{\substack{0\le i\le t\\0\le j\le n}},
 \quad Z_t(a)=H_t^{(1)}+aH_t,
 \quad Y_t(a)=H_t^{(2)}+aH_t^{(1)}.
\end{equation}
All square matrices inverted below are positive definite: their
quadratic forms are integrals of a nonzero polynomial's squared
modulus against a strictly positive density and a positive weight.
Define
\begin{equation}\label{mmnative:eq:variational}
 L_{n,t}(a)=C_t^*Z_t(a)^{-1}C_t,
 \qquad V_{n,t}(a)=D_t^*Y_t(a)^{-1}D_t,
 \qquad\mathcal D_{n,t}(a)=H_n-aL_{n,t}(a)-V_{n,t}(a).
\end{equation}
These quantities use only moments $\mu_0,\ldots,\mu_{2t+2}$.

\begin{mmnativetheorem}[Finite native Stieltjes and complementary bounds]
\label{mmnative:thm:finite}
For every $a>0$ the exact matrices
\[
 R_n(a)=\int\frac{v_nv_n^*}{x+a}\,d\lambda,
 \qquad B_n(a)=\int\frac{xv_nv_n^*}{x+a}\,d\lambda
\]
satisfy
\begin{equation}\label{mmnative:eq:bounds}
 \boxed{L_{n,t}(a)\preceq R_n(a)
       \preceq\frac{H_n-V_{n,t}(a)}a,\qquad
 V_{n,t}(a)\preceq B_n(a)\preceq H_n-aL_{n,t}(a).}
\end{equation}
The lower bounds increase with $t$, the upper bounds decrease, and
\begin{equation}\label{mmnative:eq:gap}
 \mathcal D_{n,t}(a)\succeq0
\end{equation}
is the exact computable width of the $B_n(a)$ interval and $a$ times
the exact width of the $R_n(a)$ interval.
\end{mmnativetheorem}
\begin{proof}
For a coefficient vector $c\in\mmnativeC^{n+1}$ write
$f(x)=v_n(x)^*c$.  The unique degree-at-most-$t$ minimizers in the
two weighted approximation problems are
\[
 q_0(x)=v_t(x)^*Z_t(a)^{-1}C_tc,
 \qquad q_1(x)=v_t(x)^*Y_t(a)^{-1}D_tc.
\]
Completing the square, including the complex cross terms, proves
the exact residual identities
\begin{align}
 c^*(R_n-L_{n,t})c
 &=\min_{q\in\mmnativeC[x]_{\le t}}
   \int(x+a)\left|\frac{f(x)}{x+a}-q(x)\right|^2d\lambda,
 \label{mmnative:eq:residual0}\\
 c^*(B_n-V_{n,t})c
 &=\min_{q\in\mmnativeC[x]_{\le t}}
   \int x(x+a)\left|\frac{f(x)}{x+a}-q(x)\right|^2d\lambda.
 \label{mmnative:eq:residual1}
\end{align}
For example the first integrand expands to
$\int|f|^2/(x+a)-2\Re\int\bar qf+\int(x+a)|q|^2$;
its finite quadratic minimum is $c^*R_nc-c^*C_t^*Z_t^{-1}C_tc$.
The second uses the same calculation with $xd\lambda$.
Both differences are positive semidefinite.  The exact identity
$aR_n+B_n=H_n$ gives both upper bounds.  Their widths are as stated.
Enlarging the polynomial trial space decreases each minimum, proving
all the monotonicity assertions.  No compressed multiplication matrix
was equated to the original operator.
\end{proof}

There is an exact finite Stieltjes representation of the first lower
bound.  Set $J_t=H_t^{-1/2}H_t^{(1)}H_t^{-1/2}>0$ and let
$J_t=\sum_\theta\theta P_\theta$ be its spectral decomposition.  Then
\begin{equation}\label{mmnative:eq:blockstieltjes}
 L_{n,t}(a)=\sum_\theta\frac{W_\theta}{a+\theta},\qquad
 W_\theta=C_t^*H_t^{-1/2}P_\theta H_t^{-1/2}C_t\succeq0,
 \quad\sum_\theta W_\theta=H_n.
\end{equation}
To verify the last equality, $C_t$ consists of the first $n+1$
columns of $H_t$, so $C_t^*H_t^{-1}C_t=H_n$.  This gives positive
finite nodes and block weights retaining the full moment mass.  It is
a bound produced by a finite moment problem, not a replacement measure
for any of the native integrals.  For $n=0$ the scalar weights sum to
$\mu_0$, not one.

The human block Gauss/Gauss--Radau context is the work of J\"orn
Zimmerling, Vladimir Druskin and Valeria Simoncini \cite{mmnative:ZDS},
especially its original-source Theorem \texttt{lem01} and Proposition
\texttt{prop:03}.  That proposition states a finite-matrix result
under positive block-Lanczos off-diagonal Gram factors.  None of those
nondeflation hypotheses is silently imported here: the preceding
proofs use the actual unbounded native measure and the explicit
degree-$t$ polynomial trial space of dimension $t+1$.  The claim is
proved by its variational residuals, not by applying a finite-matrix
theorem to an unbounded multiplication operator.  No scalar
potential-theory rate, extrapolation assertion, or empirical
complex-shift sign from that paper is used.

\section{A computable error rate and convergence at fixed poles}

\begin{mmnativetheorem}[Explicit pole-rate certificate]\label{mmnative:thm:rate}
Let $t=n+d$ with $d\ge0$.  The exact width in \eqref{mmnative:eq:gap} obeys
\begin{equation}\label{mmnative:eq:rate}
 \boxed{0\preceq\mathcal D_{n,n+d}(a)
 \preceq a^{-2d-2}H_n^{(2d+2)}.}
\end{equation}
The moment order on the right is still at most $2t+2$.
Thus the interval width for $B_n(4j^2)$ is bounded by
$4^{-2d-2}j^{-4d-4}H_n^{(2d+2)}$, and the Stieltjes interval has
the additional factor $1/(4j^2)$.
\end{mmnativetheorem}
\begin{proof}
Use the same admissible polynomial in both residual identities:
\[
 q(x)=\frac{f(x)}a\sum_{r=0}^d(-x/a)^r,
 \qquad\frac{f(x)}{x+a}-q(x)
 =\frac{(-1)^{d+1}x^{d+1}f(x)}{a^{d+1}(x+a)}.
\]
Its degree is at most $n+d$.  The width is exactly
$a(R_n-L_{n,t})+(B_n-V_{n,t})$.  Testing this same $q$ in both
minima makes the two weights add to $(x+a)^2$, so its total error
is precisely
$a^{-2d-2}\int x^{2d+2}|f|^2d\lambda$.  This proves the matrix
inequality, including its constant and exponent.
\end{proof}

This is a rate in the pole $a$, with an exact finite certificate at
each auxiliary degree.  It is not a claim that the displayed bound
decreases with $d$ at fixed $a$: native high moments may grow rapidly.
For an explicit numerical overestimate, the original envelope gives
\begin{equation}\label{mmnative:eq:momentcap}
 \int x^r d\nu_\epsilon
 =\zeta_{2(r+\epsilon)}
 \le\frac{2A_kC_\Gamma[2(r+\epsilon)]!}
 {\alpha^{2(r+\epsilon)+1}},\qquad
 \int x^r d\widetilde\nu_\epsilon
 \le\alpha\frac{2A_kC_\Gamma[2(r+\epsilon)]!}
 {\alpha^{2(r+\epsilon)+1}}.
\end{equation}
Indeed integrate the upper bound in \eqref{mmnative:eq:envelope}, dropping
only $(1+|y|)^{-1/2}\le1$ in this upper estimate.  Hence
$H_n^{(r)}\preceq\mmnativetr(H_n^{(r)})I$ and
$\mmnativetr(H_n^{(r)})=\sum_{i=0}^n\mu_{2i+r}$ give a fully explicit
constant in \eqref{mmnative:eq:rate}.  The measure itself has not been altered.

\begin{mmnativelemma}[Polynomial density with the original tail]\label{mmnative:lem:density}
For either native parity measure or either tilted parity measure,
polynomials are dense in $L^2(w(x)d\lambda)$ for every nonzero
nonnegative polynomial weight $w$.
\end{mmnativelemma}
\begin{proof}
The symmetric lift of $d\lambda=\ell(x)dx$ under $x=y^2$ is
$|y|\ell(y^2)dy$.  Its masses and integrals agree under this map;
for the present measures its density is
$|y|^{2\epsilon}m_k(y)s(y^2)^e$, where $e=0$ or $1$.
After multiplying by $w(y^2)$ it has a finite exponential moment
$\int e^{b|y|}d\lambda_w^{\rm lift}<\infty$ for every $b<\alpha$,
by \eqref{mmnative:eq:envelope} and $s\le\alpha$.
If $f\in L^2(wd\lambda)$ is orthogonal to all $x^r$, its even
lift $f(y^2)$ times this symmetric measure has all moments zero:
the even ones are the given orthogonalities and the odd ones vanish
by symmetry.  Cauchy--Schwarz gives finite absolute exponential
moments for that complex measure in a smaller strip, for example
$|\Re z|<\alpha/4$.  Its bilateral transform is analytic there,
and all its derivatives at zero vanish.  The identity theorem makes
the transform zero throughout the strip.  Fourier uniqueness for a
finite complex measure implies $f=0$ almost everywhere.  The
orthogonal complement of the polynomials is therefore zero.
\end{proof}

\begin{mmnativeproposition}[Certified terminating refinement]\label{mmnative:prop:converge}
For fixed $n$ and fixed $a>0$, both endpoints in \eqref{mmnative:eq:bounds}
converge monotonically to their actual integral, and
$\mathcal D_{n,t}(a)\to0$ in matrix norm as $t\to\infty$.
\end{mmnativeproposition}
\begin{proof}
Apply Lemma \ref{mmnative:lem:density} with $w=x+a$ and $w=x(x+a)$.
The targets $f/(x+a)$ have finite norm in both spaces by the finite
native moments.  Each residual minimum in
\eqref{mmnative:eq:residual0}--\eqref{mmnative:eq:residual1} tends to zero.  Apply this
to the finitely many coordinate vectors: the trace of the positive
gap tends to zero, and its norm is at most its trace.
\end{proof}
For a finite list of poles this gives simultaneous termination of any
specified positive gap tolerance, provided the needed native moments
are available as exact or certified inputs.  No uniform a priori rate
in the auxiliary degree, numerical conditioning assertion, or cost of
obtaining those moments is asserted.

\section{A finite original-moment route to the tilted integrals}

The preceding construction uses moments of the actual tilt
$s\nu_\epsilon$.  These are still unevaluated native integrals unless
supplied.  They must not be confused with moments of the raw Pollaczek
weight.  This section gives an alternative using only untilted native
moments, without requiring any exact tilted moment as input.

Take $d\lambda=d\nu_\epsilon$, let $b_r=(2r+1)^2$, and construct
the two Stieltjes endpoints from Theorem \ref{mmnative:thm:finite}:
\[
 R_{n,t}^-(a)=L_{n,t}(a),\qquad
 R_{n,t}^+(a)=(H_n-V_{n,t}(a))/a.
\]
For the second type of integral define
\[
 T_n(a,b)=\int\frac{xv_nv_n^*}{(x+a)(x+b)}d\nu_\epsilon,
 \qquad a,b>0.
\]
Put
\begin{equation}\label{mmnative:eq:quadraticlower}
 Z_t^{(2)}(a,b)=H_t^{(3)}+(a+b)H_t^{(2)}+abH_t^{(1)},
 \qquad T_{n,t}^-(a,b)=D_t^*[Z_t^{(2)}(a,b)]^{-1}D_t.
\end{equation}
This uses moments through degree $2t+3$, is positive definite, and
is a lower bound for $T_n(a,b)$.  The full residual identity is
\begin{equation}\label{mmnative:eq:quadres}
 c^*(T_n-T_{n,t}^-)c
 =\min_{q\in\mmnativeC[x]_{\le t}}
 \int x(x+a)(x+b)
 \left|\frac{v_n^*c}{(x+a)(x+b)}-q\right|^2d\nu_\epsilon.
\end{equation}
This follows by the same complete-square expansion as before, since
the cross term is $\int x\bar q(v_n^*c)d\nu_\epsilon$.
Lemma \ref{mmnative:lem:density} with the displayed polynomial weight proves
monotone convergence of this lower bound.

At distinct poles the exact signed partial fraction identity is
\begin{equation}\label{mmnative:eq:partialfraction}
 T_n(a,b)=\frac{bR_n(b)-aR_n(a)}{b-a}.
\end{equation}
Retaining its sign gives the upper endpoint
\begin{equation}\label{mmnative:eq:quadraticupper}
 T_{n,t}^+(a,b)=
 \begin{cases}
 [bR_{n,t}^+(b)-aR_{n,t}^-(a)]/(b-a),&b>a,\\[1mm]
 [aR_{n,t}^+(a)-bR_{n,t}^-(b)]/(a-b),&a>b.
 \end{cases}
\end{equation}
It bounds $T_n$ from above, is positive definite, and decreases
to $T_n$ as $t$ grows.  The combined interval width
$T_{n,t}^+-T_{n,t}^-\succeq0$ is an entirely finite moment
certificate.  In the actual application $a=4j^2$ and
$b=(2r+1)^2$, so the poles are always distinct.  No coincident-pole
limit or division by an unknown sign is needed.

For an integer $R\ge1$ put
\begin{equation}\label{mmnative:eq:theta}
 \theta_R=\frac4\pi\sum_{r=R}^\infty b_r^{-1}
 =\alpha-\frac4\pi\sum_{r=0}^{R-1}b_r^{-1},\qquad
 0<\theta_R\le\frac{2}{\pi(2R-1)}.
\end{equation}
The equality follows from \eqref{mmnative:eq:s} at zero.  For the inequality,
the decreasing function $(2u+1)^{-2}$ gives
$\sum_{r=R}^\infty(2r+1)^{-2}\le\int_{R-1}^\infty(2u+1)^{-2}du$.
All quantities are unscaled original poles.  Positive summation of
\eqref{mmnative:eq:s}, using $1/(x+b_r)\le1/b_r$, now proves
\begin{equation}\label{mmnative:eq:tiltKbounds}
 \begin{gathered}
 K_n:=\int v_nv_n^*s\,d\nu_\epsilon,\\
 K_{n,t,R}^-:=\frac4\pi\sum_{r<R}R_{n,t}^-(b_r)
 \preceq K_n\preceq
 \frac4\pi\sum_{r<R}R_{n,t}^+(b_r)+\theta_RH_n
 =:K_{n,t,R}^+.
 \end{gathered}
\end{equation}
For precisely the unresolved matrix \eqref{mmnative:eq:target},
\begin{equation}\label{mmnative:eq:tiltBbounds}
 \boxed{\frac4\pi\sum_{r<R}T_{n,t}^-(a,b_r)
 \preceq B_{\epsilon,n}(a)\preceq
 \frac4\pi\sum_{r<R}T_{n,t}^+(a,b_r)+\theta_RH_n.}
\end{equation}
Indeed the omitted integrand is at most
$\theta_R x/(x+a)v_nv_n^*$ in every quadratic form, hence at most
$\theta_RH_n$ after integration.  Alternatively it is bounded by
$\theta_RH_n^{(1)}/a$; either established upper matrix may be used,
without taking an unsupported entrywise minimum of matrices.
The lower endpoints are strictly positive, increase with $t$ and $R$,
and converge to the actual matrices.  First take $R$ large enough to
control the explicit tail, and then refine the finite list of poles
using Proposition \ref{mmnative:prop:converge} and \eqref{mmnative:eq:quadres}.  The
upper bounds converge by this same two-stage choice.  No monotonicity
in $R$ of the mixed-pole upper formula is needed.  Thus finitely many
untilted native moments suffice for any certified positive tolerance.

There is an important exact boundary to this assertion: if all
untilted native moments through degree $2n$ are already known, the
original parity Gram $H_n$ itself is directly known.  One need not
reconstruct that Gram by an auxiliary resolvent series.  The additional
formulas above certify the genuinely different $s$-weighted and
rational-weighted integrals; they do not manufacture moment values or
claim an unnecessary computational advantage.

\section{Both blocks and the complete native source certificate}

Let $N\ge q-1$ be the original polynomial cutoff, with
$n_0=\lfloor N/2\rfloor$, $n_1=\lfloor(N-1)/2\rfloor$.  Denote
the original direct-sum source Gram by
\[
 H=\mmnativediag\left(\int v_{n_0}v_{n_0}^*d\nu_0,
                  \int v_{n_1}v_{n_1}^*d\nu_1\right).
\]
For the first implementation route, the moment input is the actual
tilted sequence $\widetilde\mu_{\epsilon,r}=\int x^rsd\nu_\epsilon$;
its exact compatibility is
$\widetilde\mu_{1,r}=\widetilde\mu_{0,r+1}$.  Define the block
matrices $K=\mmnativediag(K_0,K_1)$ and
$K_x=\mmnativediag(K_{x,0},K_{x,1})$ from its unshifted and once-shifted
moments.  For every $j=1,\ldots,J$ apply \eqref{mmnative:eq:bounds} in
both blocks with $a_j=4j^2$, choosing any auxiliary degrees
$t_{\epsilon,j}\ge n_\epsilon$.  Call the resulting direct sums
$B_j^-,B_j^+$ and $\mathcal D_j=B_j^+-B_j^-$.  Set
\begin{equation}\label{mmnative:eq:Gpm}
 \begin{gathered}
 G^-_{N,J,\mathbf t}=\frac2\pi K+\frac4\pi\sum_{j=1}^JB_j^-,\\
 G^+_{N,J,\mathbf t}=\frac2\pi K+\frac4\pi\sum_{j=1}^JB_j^+
                       +\frac1{\pi J}K_x,\\
 E_{N,J,\mathbf t}=\frac4\pi\sum_{j=1}^J\mathcal D_j,
 \qquad\Delta G=G^+-G^-=E_{N,J,\mathbf t}+\frac1{\pi J}K_x.
 \end{gathered}
\end{equation}
The original reciprocal series \eqref{mmnative:eq:s} gives
\begin{equation}\label{mmnative:eq:sourcebounds}
 \boxed{0<G^-\preceq H\preceq G^+,\qquad
 0\preceq H-G^-\preceq\Delta G.}
\end{equation}
To check the outer tail including its constants, sum its positive
terms and use $x/(x+4j^2)\le x/(4j^2)$ and
$\sum_{j>J}j^{-2}\le1/J$.  The result is exactly
$K_x/(\pi J)$, not a raw Pollaczek moment block.  All inner errors
are the computed $\mathcal D_j$; both parity blocks remain present.
First choosing $J$ and then the finitely many auxiliary degrees gives
arbitrarily tight absolute brackets by the preceding convergence proofs.

Alternatively use \eqref{mmnative:eq:tiltKbounds} and
\eqref{mmnative:eq:tiltBbounds} to give the $K$ and $B_j$ lower and upper
matrices from untilted native moments.  Add the actual native outer
tail bound
\begin{equation}\label{mmnative:eq:untiltedouter}
 H-H^{[J]}\preceq\frac1{2J}X_H,\qquad
 X_H=\mmnativediag\left(\int xv_{n_0}v_{n_0}^*d\nu_0,
                    \int xv_{n_1}v_{n_1}^*d\nu_1\right),
\end{equation}
which follows from $s\le\alpha=\pi/2$.  This produces the same
type of positive bracket \eqref{mmnative:eq:sourcebounds} from a finite
untilted moment sequence only.  Each finite tail and each moment gap
must be included; no exact tilted integral is assumed in this route.

For either route the fully computable relative certificate is
\begin{equation}\label{mmnative:eq:eta}
 \eta_N=\lambda_{\max}((G^-)^{-1/2}\Delta G(G^-)^{-1/2})\ge0,
 \qquad G^-\preceq H\preceq G^+\preceq(1+\eta_N)G^-.
\end{equation}
This auxiliary conjugation computes a number; it does not change the
native source coordinates or mass.

\section{Exact complex-source and observation transport}

Retain the original coefficient map from \cite{mmnative:PP}:
\begin{equation}\label{mmnative:eq:transport}
 T_{rj}=\binom jr c^{j-r}i^r\ (r\le j),\quad
 (T^{-1})_{rj}=\binom jr(-c)^{j-r}i^{-j}\ (r\le j),\quad
 L_N=\Pi_NT,
\end{equation}
with other entries zero and $\Pi_N$ the even-first, odd-second
permutation.  It sends the original $S$-coefficients to $(E,O)$ in
\eqref{mmnative:eq:paritynorm}.  Thus
\begin{equation}\label{mmnative:eq:Sbounds}
 H_S=L_N^*HL_N,\qquad
 L_N^*G^-L_N\preceq H_S\preceq L_N^*G^+L_N.
\end{equation}
No powers of $i$ or determinant signs disappear:
\[
 \det T=i^{N(N+1)/2},\qquad
 \det L_N=(\det\Pi_N)i^{N(N+1)/2}.
\]
In particular all absolute source determinant factors are one in
this congruence, while its complex determinant phase is retained.

The original physical-unit coefficient-to-class map is $J_{S,N}$,
and the original onto observation is $\mathcal A_S=\Lambda J_{S,N}$.
In parity coordinates it is exactly
\begin{equation}\label{mmnative:eq:A}
 A=\mathcal A_SL_N^{-1}=\Lambda J_{S,N}L_N^{-1}.
\end{equation}
Every original relation and every observation-kernel direction
therefore remains in the same fibre.  For any positive source matrix
$D$, completing the square gives
\begin{equation}\label{mmnative:eq:minimum}
 Q_A(D)=(AD^{-1}A^*)^{-1},\qquad
 \min_{Az=u}z^*Dz=u^*Q_A(D)u,\quad
 z_{\min}=D^{-1}A^*Q_A(D)u.
\end{equation}
Indeed $Az_{\min}=u$ and $z_{\min}$ is $D$-orthogonal to every
vector in $\ker A$.  Restriction to a specified original subspace
uses its same inclusion in $L_N$ coordinates and the same proof on
the restricted fibre; the full pair of polynomial spaces is not
substituted for a constrained source.

Consequently \eqref{mmnative:eq:sourcebounds} and \eqref{mmnative:eq:eta} imply
\begin{equation}\label{mmnative:eq:Qbounds}
 \boxed{Q_N^-:=Q_A(G^-)\preceq Q_N:=Q_A(H)
 \preceq Q_N^+:=Q_A(G^+)\preceq(1+\eta_N)Q_N^-.}
\end{equation}
This follows by minimizing each inequality on the identical affine
fibre, not by an operator-norm replacement for $A$.  The exact lift
back to the original source coefficients is
$L_N^{-1}D^{-1}A^*Q_A(D)u$.

For any fixed original observed vectors $u_i,u_j$, the entire
positive remainder $Q_N-Q_N^-$ is retained and satisfies
\begin{equation}\label{mmnative:eq:crosserror}
 |u_i^*(Q_N-Q_N^-)u_j|^2
 \le (u_i^*\Delta Q_Nu_i)(u_j^*\Delta Q_Nu_j),\qquad
 \Delta Q_N=Q_N^+-Q_N^-.
\end{equation}
This is Cauchy--Schwarz in the positive remainder, followed by its
upper Loewner bound.  For the actual arithmetic classes put
$u_i=\Lambda a_i$ and $w=\Lambda v$, with the original
$a_1=b_N,a_2=b_{N+1}$, terminal class, and scale $\omega_N$.
If $B_i=u_i^*Q_Nw$, $B_i^-=u_i^*Q_N^-w$ and
$e_i=\sqrt{(u_i^*\Delta Q_Nu_i)(w^*\Delta Q_Nw)}$, then
\begin{equation}\label{mmnative:eq:signedcurrent}
 \left|\frac{2\Re(\bar B_1B_2-\overline{B_1^-}B_2^-)}{\omega_N}\right|
 \le\frac2{\omega_N}
 (|B_1^-|e_2+|B_2^-|e_1+e_1e_2).
\end{equation}
Expanding the exact product retains the two linear terms and the
quadratic error.  This is a certificate for approximation of that
signed current, not an assignment of its phase or sign.

\section{The full signed four-endpoint enclosure and finite tolerances}

For the original observation image rank $d_N$, let
\begin{equation}\label{mmnative:eq:deterror}
 e_N=\log\det Q_N^+-\log\det Q_N^-.
\end{equation}
These are determinants in the unchanged image frame.  The eigenvalue
bound from \eqref{mmnative:eq:Qbounds} gives
\begin{equation}\label{mmnative:eq:detcap}
 0\le\log\det Q_N-\log\det Q_N^-
 \le e_N\le d_N\log(1+\eta_N).
\end{equation}
The empty determinant is one and its contribution is zero.
At the four original cutoffs $N_1,N_2,N_3,N_4$, retain the signs
\[
 F=\log\frac{\det Q_{N_1}\det Q_{N_2}}
                    {\det Q_{N_3}\det Q_{N_4}},\qquad
 F^-=\log\frac{\det Q_{N_1}^-\det Q_{N_2}^-}
                     {\det Q_{N_3}^-\det Q_{N_4}^-}.
\]
For the original complete window these cutoffs are
$(q-1,q,2q-1,2q)$; no new endpoints are selected.  Independent
truncation lengths and moment degrees are permitted at each cutoff.
The complete signed bound is
\begin{equation}\label{mmnative:eq:four}
 \boxed{-e_{N_3}-e_{N_4}\le F-F^-
                       \le e_{N_1}+e_{N_2}.}
\end{equation}
Equivalently its lower endpoint uses $Q^-$ in the two numerator
factors and $Q^+$ in the two denominator factors; its upper endpoint
uses the reverse choices.  This keeps the endpoint signs rather than
charging all four positive errors to both sides.

Here is a finite tolerance selection for the tilted-moment route that
combines the root continuation's outer bound with the new inner
certificates.  Retain its actual constants \cite{mmnative:TR}:
\begin{equation}\label{mmnative:eq:Lsharp}
 \begin{gathered}
 \Delta_N=[1+4(N+M)^2]^M,\quad r_k=\lfloor k/3\rfloor,\\
 L_N^\sharp=\min\left\{4A_k\Delta_N(N+1)^2/a_k,
                  C_{\rm mult}(h)^2(2N+r_k)^2\right\},\\
 d_*=\alpha-\log(5/3),\quad C_\sigma=4/(\sqrt{15}\sqrt{2\pi}),\\
 H_h=\frac{64C_\sigma(1+4M^2)^M}{d_*^2}
                   \max_{b=3,4,5}\frac{M_\alpha^b}{a_b},\\
 T_h=\frac2{d_*}\{\log16+2M+3+\max(0,\log H_h)\},
 \quad C_{\rm mult}(h)=\sqrt{T_h^2+1}.
 \end{gathered}
\end{equation}
Here $a_b$ is \eqref{mmnative:eq:constants}'s original convolution constant
at order $b$.  TR2--4 proves
\begin{equation}\label{mmnative:eq:outerrelative}
 (1-\epsilon_{N,J})H\preceq H^{[J]}\preceq H,
 \qquad\epsilon_{N,J}=L_N^\sharp/(2J)<1.
\end{equation}
Since $G^-\preceq H^{[J]}\preceq G^-+E_{N,J,\mathbf t}$,
the finite moment test
\begin{equation}\label{mmnative:eq:innertest}
 E_{N,J,\mathbf t}\preceq\delta_NG^-
\end{equation}
implies
\begin{equation}\label{mmnative:eq:jointrelative}
 G^-\preceq H\preceq
 \frac{1+\delta_N}{1-\epsilon_{N,J}}G^-.
\end{equation}
The same factor applies to the observation minimum.  No integral
appears in \eqref{mmnative:eq:innertest} once the finite tilted moments are
given.  For fixed $J$, Proposition \ref{mmnative:prop:converge} makes its left
side tend to zero, while $G^-\succeq(2/\pi)K>0$, so the test is
eventually satisfied for every $\delta_N>0$.

For a prescribed four-endpoint tolerance $t_*>0$ and $d_N>0$, choose
the actual integers and positive inner tolerance
\begin{equation}\label{mmnative:eq:tolerance}
 J_N\ge\left\lceil L_N^\sharp\max\{1,4d_N/t_*\}\right\rceil,
 \qquad\delta_N=e^{t_*/(4d_N)}-1,
\end{equation}
and refine the finitely many moment degrees until \eqref{mmnative:eq:innertest}
holds.  Then $\epsilon_{N,J}\le1/2$ and
$\epsilon_{N,J}\le t_*/(8d_N)$, so
\[
 d_N\log\frac{1+\delta_N}{1-\epsilon_{N,J}}
 \le t_*/4+2d_N\epsilon_{N,J}\le t_*/2.
\]
Using these four certified individual bounds in \eqref{mmnative:eq:four}
proves $|F-F^-|\le t_*$.  Zero-rank endpoints need no such inner
test because their determinant error is zero.  At $N=O(q)$ and
$d_N\le q$ the number of outer resolvent terms is still
$O_h(q^3/t_*)$ for $0<t_*\le1$, as in \cite{mmnative:TR}.  The newly closed
inner step uses finitely many native moments, with a computable
stopping test.  No unproved bound on the auxiliary moment depth or
the cost of certifying those moments is attached to that outer rate.

\section{The same original conductor and its full kernel}

On the already established simple-quartet conductor domain, retain
the exact factorization from \cite{mmnative:EQ,mmnative:TR},
\[
 C=\bar C\Lambda,
 \qquad T_N=(\bar C Q_N^{-1}\bar C^*)^{-1},
 \qquad R_N=Q_N^{-1}\bar C^*T_N,
\]
where $\bar C$ is onto its original conductor image.  Its full
kernel is not removed.  Define
\begin{equation}\label{mmnative:eq:conductor}
 T_N^\pm=(\bar C(Q_N^\pm)^{-1}\bar C^*)^{-1}
 =\bigl[(\bar CA)(G^\pm)^{-1}(\bar CA)^*\bigr]^{-1}.
\end{equation}
Both equalities follow by substitution; the second uses the same
native source matrix and the composite observation, not a newly
chosen metric.  The identical-fibre minimum gives
\[
 T_N^-\preceq T_N\preceq T_N^+
 \preceq(1+\eta_N)T_N^-.
\]
The same assertion holds with the factor in \eqref{mmnative:eq:jointrelative}
when the finite inner test is used.  Every kernel lift is present in
the minimizing fibre $\bar Cu=w$.

Keep the original cubic inclusion $I_W$ and coordinate matrix $\Psi$,
and put
\[
 \mathcal G_N^\pm=\Psi^*I_W^*T_N^\pm I_W\Psi,
 \quad\mathcal G_N=\Psi^*I_W^*T_NI_W\Psi.
\]
For the original metric-independent signed-state map
$B_{\mathcal J}=O(I-D_{\mathcal J})O^{-1}E$ of EQ16--19, congruence
gives, for every original label vector $\alpha_0$,
\begin{equation}\label{mmnative:eq:defect}
 \alpha_0^*B_{\mathcal J}^*\mathcal G_N^-B_{\mathcal J}\alpha_0
 \le\alpha_0^*B_{\mathcal J}^*\mathcal G_NB_{\mathcal J}\alpha_0
 \le\alpha_0^*B_{\mathcal J}^*\mathcal G_N^+B_{\mathcal J}\alpha_0.
\end{equation}
The sign-flip set $\mathcal J$ is distinct from the positive integer
truncation length $J$.  The unchanged terminal-class residual in
EQ20--21 is not assigned zero by this positive energy bound.  For
conductor determinants or a retained positive definite restriction,
the identical proof of \eqref{mmnative:eq:four} uses its actual four endpoint
ranks and signs.

\section{What is certified, and exact finite tests}

The finite inputs are actual native moments, not guessed asymptotics.
With tilted moments through the stated degrees, all matrices in
\eqref{mmnative:eq:Gpm} are finite matrix expressions.  With untilted moments
only, \eqref{mmnative:eq:tiltKbounds}--\eqref{mmnative:eq:tiltBbounds} provide a second
finite route.  The original moments themselves remain unevaluated
unless supplied.  If their numerical enclosures are used instead of
exact values, every moment-derived matrix and inverse must be enclosed
with outward error bounds and its positive-definite lower margin
certified; inserting midpoint moments and treating their inverse as
exact is not a certificate.

The companion \texttt{verify\_moment\_enclosures.py} checks the
universal formulas in exact rational and Gaussian-rational arithmetic
on finite positive atomic moment systems.  Those systems are test
instances for the proved universal matrix identities, explicitly not
substitutes for $\nu_0,\nu_1$ or their tilts.  It checks both parity
blocks, both complementary moment bounds and residuals, the exact
pole-rate matrix inequality, the signed partial fractions, the
quadratic-denominator bounds, original complex coefficient congruence,
the unchanged observation and conductor minima, and signed endpoint
error inequalities.  Symbolic checks also verify the rational
identities, with all signs and powers retained.  The all-degree native
proofs above, rather than these finite examples, establish the claims
for the original arithmetic measure.

\begin{thebibliography}{9}
\bibitem{mmnative:ZDS}
J\"orn Zimmerling, Vladimir Druskin and Valeria Simoncini,
\emph{Monotonicity, bounds and acceleration of Block Gauss and
Gauss--Radau quadrature for computing $B^T\phi(A)B$},
arXiv:2407.21505v3, original author \LaTeX\ title and source
\path{GaussRadauRevised.tex}, Sections 3--5, especially original-source
labels \texttt{lemma:T\_LDLt}, \texttt{prop:Pencil2Stieltjes},
\texttt{lem01} and \texttt{prop:03}.
\url{https://arxiv.org/abs/2407.21505v3}.
The abstract-page title uses ``extrapolation'' in place of
``acceleration''; this citation identifies the retained version-three
author source.  The integrating task read that complete source; this
independent derivation checked its author/title metadata and the full
monotonicity/bounds statements and proofs surrounding the latter two
labels.  The native unbounded-measure argument is proved above.
\bibitem{mmnative:ALM}
A. I. Aptekarev, G. L\'opez Lagomasino and A. Mart\'{\i}nez--Finkelshtein,
\emph{Strong asymptotics for the Pollaczek multiple orthogonal
polynomials ensembles}, arXiv:1410.1261v1 (2014), Introduction,
unscaled weights and their positive Stieltjes ratio.
\url{https://arxiv.org/abs/1410.1261v1}.
The original author source was read by the integrating task; no
multiple-orthogonal asymptotic is imported into the native measure.
\bibitem{mmnative:PP}
Split-Zero programme, \emph{Exact native parity transport into the
Pollaczek pair and certified finite source--observation bounds},
20 September 2026, complete source
\path{independent_parity_pollaczek/PARITY_POLLACZEK_NATIVE_RECEIVER.tex};
original parity maps and masses, positive Green-kernel series,
finite Stieltjes compression and source--observation bounds.
Read completely for this derivation.
\bibitem{mmnative:TR}
Split-Zero programme, \emph{Native resolvent truncation at growing
degree and the original conductor minimum}, 20 September 2026,
complete source \path{NATIVE_TRUNCATION_CONDUCTOR_RETURN.tex},
TR1--10.  Read completely for this derivation.  The native outer
truncation, tensor multiplication constant, full signed endpoint
calculation and exact conductor receiver are retained here.
\bibitem{mmnative:CAI}
Split-Zero programme, \emph{The complete original action relative
to its arithmetic total}, complete provider
\path{CAI_COMPLETE_PREREQUISITE.tex}, CAI16--18 and CAI21--25,
with its complete original convolution providers and original constants.
\bibitem{mmnative:HG}
Split-Zero programme, \emph{Subexponential heat cost and the original
arithmetic currents}, 19 September 2026, HG21, arbitrary fixed
primary multiplicity and original full-measure comparison;
complete provider \path{HEAT_GROWTH_AND_ORIGINAL_CURRENTS.tex}.
\bibitem{mmnative:EQ}
Split-Zero programme, \emph{The eight signed states through the
original conductor minimum}, complete companion
\path{EIGHT_STATE_ORIGINAL_MINIMUM.tex}, EQ1--21, and the independent
\path{EIGHT_STATE_QUOTIENT_REDERIVATION.tex}, IQ1--29.
The exact receiving maps used here are the equations retained and
proved in the complete root continuation \cite{mmnative:TR}; no new
conductor-domain theorem is asserted in this note.
\end{thebibliography}

\endgroup

% END COMPLETE NATIVE FINITE-MOMENT PROOFS026

\end{document}
